% PDF-derived LaTeX transcription. Original manuscript source was unavailable.
% Published PDF SHA-256: 3015fdb73f33ddac1c5b999d498504f8b5da17f247cad15ee64a4a250e386e1b
% Extracted text SHA-256: 50197f51d093134dc7fefcb99f4e452b6d5fafda647e3648dee30c09fea485bf
\documentclass[11pt]{article}
\usepackage[margin=1in]{geometry}
\usepackage{amsmath,amssymb}
\setlength{\emergencystretch}{3em}
\title{Contraction-precision sensitivity in a coupled D3Q19 rheology–particulate workflow}
\author{}
\date{}
\begin{document}
\maketitle
\noindent\textbf{PDF-derived LaTeX transcription.} This editable text follows the PDF extraction. Refer to the published PDF for figures, table layouts, and mathematical notation.

% Extracted chunk 1

Contraction-precision sensitivity in a coupled D3Q19
rheology–particulate workflow

\section*{Abstract}

Floating-point storage type alone does not determine the arithmetic of compiled lattice
Boltzmann contractions. We examine contraction-precision sensitivity in one coupled D3Q19
rheology–particulate workflow using paired saved reports for two particulate initializations at
identical physical parameters, a 1283 grid, and 1600 lattice steps. Changing only the reference
stepper’s matrix-multiplication precision context from default to highest substantially changed
its saved trajectory summaries: default-to-highest final specific kinetic-energy ratios were 9.8175
and 9.7829, and final mean-shear ratios were 3.1239 and 3.1162. A structure-of-arrays imple-
mentation and a Pallas implementation had equal recorded scalar and channel summaries; the
structure-of-arrays implementation agreed much more closely with the highest-precision-context
reference than with the default reference. Agreement was not exact: the largest remaining chan-
nel normalized root-mean-square error was 0.006783. Stage diagnostics and captured compiler
output support contraction arithmetic and cancellation in near-equilibrium population differ-
ences as an explanation, without establishing an effective numerical-force field. Both distinct
trajectory families passed the saved harness audits. This case therefore distinguishes internal
audit consistency from trajectory fidelity and identifies contraction precision as a consequential
numerical control in this particular workflow, rather than establishing a continuum solution or
a general result for lattice Boltzmann rheology.

\section*{1 Introduction and contributions}

A numerical reference can be reproducible under its own implementation contract yet remain sen-
sitive to arithmetic choices that are not apparent from array dtype. This distinction matters in
lattice Boltzmann workflows where small differences between populations and their equilibrium
values feed a shear estimate, a viscosity update, and subsequent collision steps. Internal identities
and admissibility tests can remain satisfied even when the resulting velocity and shear trajectories
differ substantially.
This paper studies one such case: a forced, isothermal D3Q19 Bhatnagar–Gross–Krook (BGK)
workflow coupled to shear-dependent viscosity and particulate transport in a branched-pipe wall
geometry. The evidence comprises four saved paired comparison reports, stage diagnostics, a single-
population perturbation control, source snapshots, and captured compiler records. The two samples
share all physical parameters and differ only in their initial particulate seed. They are not two
physical systems or a parameter sweep.
The focused contribution is threefold. First, a comparison retaining the reference stepper while
changing its matrix-multiplication precision context links contraction precision to a large change
in final bulk diagnostics. Second, alternative implementations agree more closely at the saved-
summary level with the precision-controlled reference, while retaining measurable channel errors.
Third, compiler, stage, and control records help distinguish a repeated arithmetic effect from the
1

response to one isolated initial perturbation. Together these results show why passing a consistency
audit does not by itself establish fidelity to a particular numerical trajectory. They do not establish
a new constitutive model, an independently validated continuum solution, or the prevalence of this
behavior across other lattice Boltzmann implementations.

\section*{2 Related work}

Precision sensitivity in lattice Boltzmann methods is not a new general observation. Lehmann
et al. compare 64-bit, 32-bit, and customized 16-bit formats across several systems, including
arrangements that separate reduced-precision storage from FP32 arithmetic [ 1]. Their favorable
FP32 findings are conditional on those implementations and test systems. They neither validate
the present coupled workflow nor provide a matched numerical benchmark for it.
The relevant distinction here is between stored dtype and contraction arithmetic. JAX docu-
ments controls for the default precision of matrix multiplication and convolution with 32-bit inputs,
including a context manager and per-operation settings [ 2]. Thus, float32 inputs do not uniquely
specify the arithmetic used by a compiled contraction. Contemporary API documentation explains
this distinction but does not establish the historical JAX version, backend configuration, or ex-
ecutable used for any saved report. The present contribution is the controlled comparison and
associated diagnostic evidence within one recorded workflow.

\section*{3 Scientific methods and validation protocols}

Study design and evidence scope
This is a retrospective numerical case study; no additional solver simulations were performed for
this manuscript. Four historical reports compare jax-soa, pallas, and jax-highest against
the default jax reference, and jax-soa against jax-highest. Comparisons remain paired by
particulate initialization, denoted A and B. Each trajectory uses a 1283 grid, 1600 lattice steps,
and 16 saved frames at steps 100, 200, . . . ,1600. These constitute two seed-defined initializations,
not independent physical replicates; repeated appearances in different reports do not increase the
trajectory count.
The new independent verification consists of 150 recorded-identity, saved-arithmetic, cross-report,
control-log, and static-source/compiler checks, all of which passed. These checks did not execute
the supplied solver or recompute channel errors from original target arrays. Those arrays and the
complete historical environment are absent from the captured verification. Exact saved outputs
and original records are retained in Supplementary record S1.
2

Lattice dynamics and rheological coupling
For populations fi, velocities ci, and weights wi, the static source describes the isothermal D3Q19
relations
\ensuremath{\rho} =
18\ensuremath{\sum}
i=0
fi, u =
\ensuremath{\sum}
i fici
max(\ensuremath{\rho}, 10\ensuremath{-}8) , (1)
f eq
i = \ensuremath{\rho}wi
[
1 + ci · u
c2s
+ (ci · u)2
2c4s
\ensuremath{-} |u|2
2c2s
]
, c2
s = 1
3 . (2)
The weights are 1/3 for rest, 1/18 for axial directions, and 1/36 for diagonal directions. The source
stores C and W as float32. The non-equilibrium moment estimate is
\ensuremath{\Pi}neq
\ensuremath{\alpha}\ensuremath{\beta} =
\ensuremath{\sum}
i
(fi \ensuremath{-} f eq
i )ci\ensuremath{\alpha}ci\ensuremath{\beta}, (3)
D\ensuremath{\alpha}\ensuremath{\beta} = \ensuremath{-}
\ensuremath{\Pi}neq
\ensuremath{\alpha}\ensuremath{\beta}
2\ensuremath{\rho}c2s\ensuremath{\tau}prev
, ˙\ensuremath{\gamma} =
\ensuremath{\surd}
2
\ensuremath{\sum}
\ensuremath{\alpha},\ensuremath{\beta}
D2
\ensuremath{\alpha}\ensuremath{\beta} + 10\ensuremath{-}20. (4)
Here \ensuremath{\tau}prev is the previous-step relaxation time. With particulate solid volume fraction \ensuremath{\varphi}, the
Carreau–Yasuda viscosity and concentration correction are
\ensuremath{\mu}CY = \ensuremath{\mu}\ensuremath{\infty} + (\ensuremath{\mu}0 \ensuremath{-} \ensuremath{\mu}\ensuremath{\infty})[1 + (\ensuremath{\lambda} ˙\ensuremath{\gamma})a](n\ensuremath{-}1)/a, (5)
\ensuremath{\mu}\ensuremath{\ast} = \ensuremath{\mu}CY
[
max
(
1 \ensuremath{-} \ensuremath{\varphi}
\ensuremath{\varphi}m
, 0.05
)] \ensuremath{-}[\ensuremath{\eta}]\ensuremath{\varphi}m
, (6)
\ensuremath{\mu} = min\{\ensuremath{\mu}\ensuremath{\ast}, (\ensuremath{\tau}max \ensuremath{-} 0.5)\ensuremath{\rho}c2
s\}, (7)
\ensuremath{\tau} = clip
(
0.5 + \ensuremath{\mu}
max(\ensuremath{\rho}, 10\ensuremath{-}8)c2s
, 0.50011, \ensuremath{\tau} max
)
. (8)
Here the lattice parameters are \ensuremath{\mu}0 = 0.16, \ensuremath{\mu}\ensuremath{\infty} = 0.0035, \ensuremath{\lambda} = 8.2, n = 0.2128, a = 0.64, \ensuremath{\varphi}m = 0.63,
[\ensuremath{\eta}] = 2 .0, and \ensuremath{\tau}max = 3.0. The collision update in fluid cells is
f post
i = fi \ensuremath{-} fi \ensuremath{-} f eq
i
\ensuremath{\tau} + wi
cix
c2s
Fx, F x = 6.0 \ensuremath{\times} 10\ensuremath{-}7. (9)
Streaming uses directional rolling with the supplied wall handling. Wall velocities are zeroed, wall
post-collision populations are reset to unit-density zero-velocity equilibrium, and opposite-direction
populations are selected at wall destinations. The geometry and periodic- x rolling are historical
implementation conventions, not prescribed inlet/outlet data or a canonical analytical pipe-flow
benchmark.
Both samples start from the same source-defined fluid state: unit density, a predominantly
parabolic streamwise velocity profile with a sinusoidal axial modulation, equilibrium populations,
and initial fluid relaxation time 0.82. The initial particulate field is a smooth deterministic trigono-
metric modulation, rescaled to mean fluid solid fraction 0.08 and clipped to its prescribed bounds.
Its seed alone differs between A and B.
Particulate transport and declared units
The historical particulate update uses centered face fluxes. Overbars denote separate arithmetic
averages of concentration and particle velocity across the face; the subscript + denotes the neigh-
3

boring cell in the positive \ensuremath{\alpha} direction. For an open fluid–fluid face,
J\ensuremath{\alpha},+ = \ensuremath{\varphi} vp\ensuremath{\alpha} \ensuremath{-} Dp(\ensuremath{\varphi}+ \ensuremath{-} \ensuremath{\varphi}) \ensuremath{-} Km\ensuremath{\varphi}( ˙\ensuremath{\gamma}+ \ensuremath{-} ˙\ensuremath{\gamma}), (10)
with Dp = 2 .0 \ensuremath{\times} 10\ensuremath{-}3 and Km = 2 .0 \ensuremath{\times} 10\ensuremath{-}4 in lattice units; closed faces have zero flux. Particle
velocity includes hindered settling,
vp = u \ensuremath{-} vs max(1 \ensuremath{-} \ensuremath{\varphi}, 0)4.65ez, v s = 1.25 \ensuremath{\times} 10\ensuremath{-}4. (11)
The update is
\ensuremath{\varphi}n+1 = clip
(
\ensuremath{\varphi}n \ensuremath{-}
\ensuremath{\sum}
\ensuremath{\alpha}
[J\ensuremath{\alpha},+ \ensuremath{-} J\ensuremath{\alpha},\ensuremath{-}], 0, 0.95\ensuremath{\varphi}m
)
, (12)
followed by zeroing wall values. Face differences telescope before clipping; clipping prevents a claim
of exact conservation. The configured particle diameter is 5.0 \ensuremath{\times} 10\ensuremath{-}5 m; settling in these reports is
explicitly prescribed in lattice units.
The historical declared conversion is \ensuremath{\Delta}x = 10 \ensuremath{-}3 m and \ensuremath{\Delta}t = 10 \ensuremath{-}3 s, with density scale 1000
kg/m3, velocity scale 1 m/s, and dynamic-viscosity scale 1 Pa s. Shear and strain are divided by
\ensuremath{\Delta}t; pressure is the isothermal gauge quantity c2
s(\ensuremath{\rho} \ensuremath{-} \ensuremath{\rho}ref ) times the pressure scale, with a framewise
fluid-density reference. Extra stress is 2\ensuremath{\mu}D. This mapping is declared rather than experimentally
calibrated rheology. X64 is enabled for derived-field conversion, which uses float64 elementwise
arithmetic before packing float32 targets; it does not imply float64 lattice dynamics. State target
fields are zero in walls.
Precision intervention and comparison metrics
The default reference delegates to the batched JAX stepper. The jax-highest variant wraps that
same stepper in jax.default\_matmul\_precision("highest"); it changes contraction precision,
not the constitutive equations, and is not an FP64 truth solution. The jax-soa implementation
expresses direction sums through channel arrays and explicit additions/subtractions. Pallas pro-
vides another candidate implementation, whose role here is evaluated through its saved comparisons
rather than assumed independence of every code path.
For channel q, let cq and rq denote candidate and reference target values, respectively. The
original harness computes
nRMSEq =
\ensuremath{\surd}
\ensuremath{\langle}(cq \ensuremath{-} rq)2\ensuremath{\rangle}
max(
\ensuremath{\surd}
\ensuremath{\langle}r2q \ensuremath{\rangle}, 10\ensuremath{-}300)
. (13)
The averages pool all 16 frames and all lattice cells, including wall zeros, separately for each paired
sample. Arithmetic is float64 on the saved float32 target arrays. By contrast, final bulk values use
only fluid cells at step 1600. In particular,
Ef =
\ensuremath{\langle} |u|2
2
\ensuremath{\rangle}
fluid,1600
(14)
is a specific kinetic-energy diagnostic, not integrated or density-weighted energy. Grid cells and
repeated outputs are not independent replicates, and no confidence intervals are inferred from them.
The saved harness audits test internal properties rather than external accuracy. Particulate
mass drift is the maximum absolute fractional change in the fluid-cell solid-fraction sum relative
4

to the first saved frame (step 100), with that baseline sum floored at 10\ensuremath{-}12. It does not bound
changes during the unsaved interval from step zero to step 100. The historical snapshot includes
18 named checks, including mass drift at most 5\%, stress-identity relative error at most 2 \ensuremath{\times} 10\ensuremath{-}5
on active cells, absolute framewise gauge-pressure mean at most 10\ensuremath{-}3 Pa, Mach number below 0.3,
and relaxation-time bounds. Report prose elsewhere mentions 32 checks; this coverage gap is not
resolved by reconstructing a 32-check suite or by substituting current audit settings.

\section*{4 Experimental results}

Reference-precision sensitivity
Changing the reference precision context produces a large separation in final velocity-derived and
shear diagnostics. Table 1 reports values from the paired default-versus-highest comparison; dis-
played decimals are rounded from the saved values, not estimates with uncertainty.
Table 1: Final fluid-cell diagnostics under the declared mapping. Ef has units m 2/s2 and mean
shear has units s \ensuremath{-}1. Ratios are default divided by highest.
Quantity Precision context A B
Ef Default 3.6524636471 \ensuremath{\times} 10\ensuremath{-}6 3.6395838379 \ensuremath{\times} 10\ensuremath{-}6
Highest 3.7203656689 \ensuremath{\times} 10\ensuremath{-}7 3.7203668057 \ensuremath{\times} 10\ensuremath{-}7
Ratio 9.8174856242 9.7828628947
Mean shear Default 0.2751940787 0.2745169997
Highest 0.0880941376 0.0880941749
Ratio 3.1238637003 3.1161765253
These are ratios of final bulk diagnostics, not tenfold or threefold nRMSE values. For example,
the pooled shear-channel nRMSE of highest against default is 0.19121103888426583 for A and
0.1930672210683576 for B. The final mean streamwise velocities are approximately 0.002290 and
0.002288 m/s under default precision, compared with 0.000733525 m/s under highest precision.
The imposed force and physical parameters are unchanged.
Agreement among alternative implementations
The SoA and Pallas reports have equal saved bulk values and equal channel-error summaries for
each sample, as well as matching recorded digest prefixes. This establishes recorded agreement,
not a new independent full-byte comparison of their absent arrays. Both differ from the default
reference. Table 2 shows selected channels and the worst channel, preserving sample pairing and
reference direction.
Agreement improves strongly when highest replaces default as reference, but does not become
bit-exact or satisfy a uniform 10\ensuremath{-}6 channel tolerance. The exact worst-channel values for SoA
versus highest are 0.006776605623877569 and 0.006782976864327476. Normal stress discrepancies
remain larger than streamwise-velocity and shear discrepancies; the available summaries do not
establish that these residual errors are harmless.
At the final frame, SoA’s relative differences from highest in Ef are 1.2987138989 \ensuremath{\times} 10\ensuremath{-}5 and
1.2452370637 \ensuremath{\times} 10\ensuremath{-}5 for A and B. Corresponding mean-shear differences are 7.6117693177 \ensuremath{\times} 10\ensuremath{-}6
5

Table 2: Pooled channel nRMSE. S denotes SoA and P denotes Pallas; their saved errors against
default are equal. The worst channel in every listed comparison is stresszz .
Candidate/reference Sample ux Pressure Shear Worst
S or P/default A 0.18671153 0.52084574 0.19121006 0.84654456
B 0.18830605 0.52775796 0.19306626 0.84592513
Highest/default A 0.18671258 0.52089450 0.19121104 0.84641184
B 0.18830708 0.52780602 0.19306722 0.84579429
S/highest A 4.48666414 \ensuremath{\times} 10\ensuremath{-}6 0.00073067 2.74615570 \ensuremath{\times} 10\ensuremath{-}5 0.00677661
B 4.47299007 \ensuremath{\times} 10\ensuremath{-}6 0.00073071 2.74775086 \ensuremath{\times} 10\ensuremath{-}5 0.00678298
and 6.6814391312\ensuremath{\times}10\ensuremath{-}6. Final fluid sums of solid fraction are equal in the paired summaries. Small
bulk differences do not contradict the larger pooled normal-stress errors: they measure different
quantities over different denominators and domains.
Stage and compiler evidence
Stage probes compare reference and SoA operations at matched input within each probe. Each
precision mode first develops its own state for 100 steps, so the default and highest probes do
not share an identical developed state across modes. Their within-probe discrepancies nevertheless
provide useful localization (Table 3).
Table 3: Selected within-probe reference–SoA nRMSE values. Each column uses its own precision-
developed state; these are not cross-mode identical-state errors.
Stage Default probe Highest probe
Density 1.16505955 \ensuremath{\times} 10\ensuremath{-}7 1.16528268 \ensuremath{\times} 10\ensuremath{-}7
ux 1.09861848 \ensuremath{\times} 10\ensuremath{-}3 9.21474029 \ensuremath{\times} 10\ensuremath{-}8
uy 6.97468972 \ensuremath{\times} 10\ensuremath{-}2 9.37988987 \ensuremath{\times} 10\ensuremath{-}8
uz 1.01001352 \ensuremath{\times} 10\ensuremath{-}1 9.37362344 \ensuremath{\times} 10\ensuremath{-}8
Equilibrium 8.61107044 \ensuremath{\times} 10\ensuremath{-}6 1.65032694 \ensuremath{\times} 10\ensuremath{-}7
Non-equilibrium populations 3.55637701 \ensuremath{\times} 10\ensuremath{-}2 6.84150191 \ensuremath{\times} 10\ensuremath{-}4
Shear 5.77717181 \ensuremath{\times} 10\ensuremath{-}4 3.04549431 \ensuremath{\times} 10\ensuremath{-}4
Streamed populations 8.01760810 \ensuremath{\times} 10\ensuremath{-}6 2.73133176 \ensuremath{\times} 10\ensuremath{-}8
Density discrepancies remain near 10\ensuremath{-}7, whereas default transverse-velocity discrepancies are
much larger relative to their small reference RMS values. The largest absolute transverse-velocity
differences in the default probe are about 7.9 \ensuremath{\times} 10\ensuremath{-}5 lattice units. Under highest precision the
velocity discrepancies approach 10\ensuremath{-}7 relative. Non-equilibrium populations remain more sensitive
than equilibrium populations, as expected when subtracting nearly equal quantities. Not every
downstream stage improves proportionally: viscosity and relaxation-time nRMSE remain of similar
magnitude between probes.
Captured optimized HLO contains contraction-related GEMM fusion, and the associated PTX
contains 152 TF32 conversion sites and 64 TF32 tensor-contraction sites. These static instruction
sites establish compiled TF32 use in that captured batch-four module; they are not dynamic exe-
cution counts or per-sample whole-environment attestations. Together with the precision-context
6

comparison, the records support a contraction-precision explanation for the observed separation.
They do not measure an effective stochastic forcing distribution or close a numerical-force budget.
Perturbation control and audit outcomes
A historical control perturbed one population at one fluid site by one positive float32 unit in the
last place, then advanced both states with the same default reference. Population, relaxation-
time, and particulate discrepancies were reported as zero at the first 100-step checkpoint and every
subsequent checkpoint through step 1600. This control does not support the proposed explanation
that an arbitrary isolated one-ULP perturbation necessarily grows beyond a 10\ensuremath{-}6 field tolerance.
It also does not establish monotone decay before step 100, global stability, absence of chaos, or a
Lyapunov exponent.
An earlier short reference–SoA comparison did fail its stated 10\ensuremath{-}6 criterion, with worst reported
nRMSE 0.015720989447758318. The later highest-reference comparison represents improved agree-
ment under a changed reference contract, not a retroactive pass of that original comparison.
All four full-length reports record successful reference and candidate audits, despite the two
distinct trajectory families. Candidate particulate relative mass drift is 1.1920928955078125 \ensuremath{\times} 10\ensuremath{-}7
under the first-saved-frame baseline defined above. Other candidate metrics include maximum Mach
number about 0.053713, relaxation times approximately 0.999–1.085, and stress-identity relative
errors below 3.6 \ensuremath{\times} 10\ensuremath{-}7. These are saved harness outcomes. They establish internal consistency as
reported by that harness, not external physical accuracy.

\section*{5 Discussion and limitations}

The strongest evidence is the precision intervention within the reference stepper: changing its
contraction context reproduces the lower-energy, lower-shear family approached by the alternative
implementations. This is more discriminating than a comparison between two layouts alone, be-
cause layout changes can alter both reduction order and compiler lowering. The SoA/Pallas saved
agreement and the stage diagnostics reinforce the interpretation without making highest an exact
solution.
The numerical mechanism is plausible and source-supported. Momentum contractions can sub-
tract contributions whose net is small; equilibrium populations depend on that momentum; and
fi \ensuremath{-} f eq
i subtracts quantities much larger than their difference. A schematic perturbation satisfies
\ensuremath{\delta}f neq
i = \ensuremath{\delta}fi \ensuremath{-} \ensuremath{\delta}f eq
i , (15)
so a small absolute equilibrium error can be a large fraction of a near-equilibrium residual. Its
second moment feeds shear, which feeds viscosity and collision relaxation. Repeated arithmetic
differences in this feedback loop need not resemble a single initial perturbation. This explains why
the one-site control cannot settle the effect of a changed contraction throughout a trajectory. The
evidence supports this mechanism but does not isolate every contraction’s contribution or establish
a measured additional force.
Audit consistency and trajectory fidelity answer different questions. Bounds, gauge pressure,
and the stress identity can hold for both flow states because they test admissibility or relations
among jointly derived quantities. In particular, a stress field derived from the same strain and
viscosity as the shear diagnostic can satisfy its algebraic identity while differing from another
7

implementation’s field. Neither audit success nor agreement among alternative implementations
substitutes for continuum validation.
Several scientific limits remain. There is no grid refinement, time-step refinement, parameter
sweep, calibrated rheological experiment, analytical benchmark, or FP64 dynamical reference. The
two closely related particulate initializations provide paired corroboration at one configuration,
not a population estimate. Pooled nRMSE includes walls and all saved times, while final bulk
diagnostics emphasize fluid cells at one time; neither reveals the spatial structure of residual errors.
The remaining normal-stress error warrants attention rather than dismissal.
Provenance also bounds the interpretation. The currently captured raw generator explicitly
disables TF32 and differs from the historical snapshot. That snapshot supports method inter-
pretation, not per-job source attestation. Earlier reports omit the reference-kernel field emitted
by the later captured validation source, and the historical environment is not immutably bound
to every report. Furthermore, the validation source bypasses its published-digest comparison for
nondefault references or non-1600-step runs. Consequently, a true reference-reproduction flag in
the highest-reference report is not evidence that it matches the default trajectory; its saved digest
prefixes differ. The 18-versus-32 audit-count discrepancy likewise remains a source/run coverage
gap. Supplementary record S1 preserves these distinctions alongside the exact records.
A prospective investigation could use identical developed states across precision modes, preserve
target arrays and complete executable provenance, and compare against higher-precision dynamics
and refinement studies. Those tests were not performed here. Operational suitability was not
evaluated.

\section*{6 Conclusion}

In this fixed D3Q19 rheology–particulate workflow, contraction precision materially changes saved
trajectory diagnostics despite unchanged physical parameters and float32 lattice storage. Default-
to-highest final specific kinetic-energy ratios are approximately 9.8 and mean-shear ratios approx-
imately 3.12 for both particulate initializations. SoA and Pallas have equal recorded summaries,
and SoA agrees much more closely with the highest-context reference than with default, although
its worst pooled channel error remains approximately 0.006783.
The controlled precision comparison, stage diagnostics, compiler evidence, and isolated perturba-
tion control support a contraction-sensitive numerical mechanism rather than a general assertion of
perturbation growth. Successful saved audits coexist with distinct flow states. The resulting lesson
is specific but important: internal consistency, reference reproduction, and trajectory fidelity must
be evaluated separately when contraction arithmetic is part of the numerical method.

\section*{7 Data and Code A vailability}

The evidence exists as private captured comparison reports, diagnostic logs, source captures, histor-
ical source snapshots, compiler records, and a saved-evidence verification recipe and result. Exact
original records and numerical outputs are preserved in the private Supplementary record S1 (ID:
S1). The captured sources include the comparison harness, precision wrapper, SoA implementa-
tion, stage and perturbation scripts, and target-conversion routines. The verification recipe checks
recorded identities, arithmetic, cross-report consistency, control records, and static source/compiler
relationships; it is not a solver-reproduction script.
8

The original target arrays are not included in this captured verification, and the original harness
states that its comparison arrays were held in memory rather than written as target stores. Com-
plete historical environment provenance is also unavailable here. Access to the captured materials is
through the existing private supplementary record; no public archive, reviewer-access entitlement,
or availability-on-request permission is established. Thus the supplied record supports inspection of
saved evidence, not independent reconstruction of the original trajectories from a complete archived
execution environment.

\section*{References}

[1] M. Lehmann, M. J. Krause, G. Amati, M. Sega, J. Harting, and S. Gekle. Accuracy and
performance of the lattice Boltzmann method with 64-bit, 32-bit, and customized 16-bit number
formats. Physical Review E 106, 015308 (2022). doi:10.1103/PhysRevE.106.015308.
[2] The JAX Authors. jax.default\_matmul\_precision. JAX API documentation. https://docs.
jax.dev/en/latest/\_autosummary/jax.default\_matmul\_precision.html. Accessed 26 Oc-
tober 2026.
9

\end{document}
